% \noindent \method includes three domains that cover a diverse range of real-world applications and challenges, but there is room for expansion to a wider range of tasks.
% Thanks to \method's integration with Gymnasium \citep{towers_gymnasium_2023} (previously OpenAI Gym) and the implementation of baselines using TF-Agents \citep{TFAgents}, adding new domains and baselines is straightforward, making it easy for researchers to contribute to the platform.
% % In partnership with the Farama Foundation\footnote{https://farama.org/about}, an organization that provides long-term maintenance to many other widely used open-source RL benchmarks such as MiniGrid \cite{MinigridMiniworld23}, D4RL \cite{fu2020d4rl}, MetaWorld \cite{yu2019meta}, MiniWob++ \cite{liu2018reinforcement}, and the Arcade Learning Environment \cite{machado18arcade}, we will ensure that \method is well-maintained and continues to evolve to meet the needs of future researchers. 

% Future work could expand \method to include multi-agent domains and tasks, reflecting real-world scenarios where autonomous agents interact with other agents and humans. Additionally, adding support for measuring system performance on customized hardware platforms would provide more precise insights into performance in target deployment environments, as current evaluations are primarily conducted on desktop and server machines. Another area of future work is further standardizing evaluations in \method, addressing potential variations due to different hardware setups, Python versions, and code implementations. These efforts will enhance reproducibility and facilitate more accurate comparisons across different research environments, further solidifying \method's role as a comprehensive benchmark for real-world autonomous agents.


% \section{Limitations and Future Work}

\method includes three domains that cover a diverse range of real-world applications and challenges, but there is room for expansion to a wider range of tasks. Thanks to \method's integration with Gymnasium \citep{towers_gymnasium_2023} (previously OpenAI Gym) and the implementation of baselines using TF-Agents \citep{TFAgents}, adding new domains and baselines is straightforward, making it easy for researchers to contribute to the platform.

Future work could expand \method to include multi-agent domains and tasks, reflecting real-world scenarios where autonomous agents interact with other agents and humans. Many real-world applications inherently involve multiple agents coordinating or competing: autonomous vehicles navigating in traffic, robot teams in warehouse settings, or trading agents in financial markets. Integrating multi-agent domains would require additional metrics to capture interaction dynamics, such as coordination efficiency, communication overhead, and emergent social behaviors. For example, extending the quadruped locomotion domain to include multiple robots collaboratively navigating complex terrain would test both individual control and collective coordination capabilities, potentially revealing new insights about algorithm robustness in social contexts.

Another area of future work is the addition of support for measuring system performance on custom hardware platforms.
This would provide more precise insights into performance in target deployment environments, as current evaluations are conducted primarily on desktop and server machines.
This extension is particularly important for edge computing applications such as robotics, where power constraints, thermal limitations, and specialized accelerators significantly influence real-world performance.
By developing standardized benchmarking procedures for specific deployment platforms such as NVIDIA Jetson\footnote{\url{https://www.nvidia.com/en-us/autonomous-machines/embedded-systems/}}, Google Coral\footnote{\url{https://coral.ai/}}, or custom FPGA implementations, \method could offer more accurate predictions of deployment performance. This would help bridge the gap between research prototypes and production systems by identifying specific optimization opportunities for the target hardware.

To further standardize evaluations in \method, future work should address potential variations due to different hardware setups, Python versions, and code implementations.
Even with our current efforts to ensure reproducibility, subtle differences in environment configurations can lead to meaningful performance variations.
Creating a centralized evaluation server, similar to the approach taken by MLPerf \cite{reddi2020mlperf}, could further standardize comparisons by running all submissions in identical environments. These enhancements would facilitate more accurate comparisons between different computing environments.

As an open-source platform, \method is designed to evolve through community contributions. Researchers can extend the benchmark in multiple ways: by adding new domains through the standard Gymnasium interface, by implementing additional baseline algorithms, or by introducing domain-specific metrics that capture other aspects of real-world performance. The repository includes detailed contribution guidelines, templates, and documentation to facilitate these extensions. This collaborative approach ensures \method remains relevant to emerging research challenges while expanding its coverage of real-world autonomous agent applications.
